\documentclass[a4paper,11pt]{article}
\usepackage[utf8]{inputenc}
\usepackage[T1]{fontenc}
\usepackage[english]{babel}
\usepackage[left=1cm,right=1cm,top=.5cm,bottom=0cm]{geometry}
\usepackage{enumitem}
\usepackage{hyperref}
\usepackage{xcolor}
\usepackage{titlesec}
\usepackage{marvosym}
\usepackage{lmodern}
\usepackage[normalem]{ulem}

\definecolor{primary}{RGB}{0, 51, 102}
\definecolor{secondary}{RGB}{80, 80, 80}

\hypersetup{colorlinks=true, linkcolor=primary, filecolor=primary, urlcolor=primary}
\titleformat{\section}{\large\bfseries\color{primary}\uppercase}{}{0em}{}[\titlerule]
\titlespacing{\section}{0pt}{8pt}{5pt}
\newcommand{\resumeItem}[1]{\item\small{#1}}

\begin{document}
\begin{center}
    {\Large \textbf{Lo\"ic Aron MBASSI EWOLO}} \\ \vspace{4pt}
    \textbf{\large Data Scientist -- Intelligence Artificielle} \\
    \textit{\small Stage INRIA Saclay (Avril--Ao\^ut 2026) -- Disponible en alternance d\`es septembre 2026 (12 mois)} \\ \vspace{4pt}
    \small
    \Pointinghand\ Cergy, France (Mobile IDF) \quad
    \Mobilefone\ (+33) 7 59 52 33 98 \quad
    {\color{primary}\Letter\ \href{mailto:loicaron.mbassiewolo@ensea.fr}{loicaron.mbassiewolo@ensea.fr}} \\ \vspace{2pt}
    {\color{primary}LinkedIn: \href{https://www.linkedin.com/in/loic-mbassi/}{linkedin.com/in/loic-mbassi}} \quad
    {\color{primary}GitHub: \href{https://github.com/Nameless0l}{github.com/Nameless0l}} \quad
    {\color{primary}Portfolio: \href{https://mbassiloic.tech}{mbassiloic.tech}}
\end{center}
\vspace{-6pt}
\noindent\small{Data Scientist IA avec expertise en NLP applique aux donnees financieres et processus de scoring de risque. Actuellement en stage INRIA Saclay (Avril 2026) sur pipeline NLP de scoring confiance/transparence, j'apporte une experience operationnelle de ML/NLP production. Double diplome ENSEA/Ecole Polytechnique Yaounde (mathematiques appliquees). Experience recherche MILA Montreal sur generalisation reseau neuronaux et deep learning. Passionnee par l'innovation IA chez BPCE et transformation numerique bancaire.}

\section{Formation}
\begin{itemize}[leftmargin=*, label=\tiny$\bullet$, nosep]
    \item \textbf{ENSEA -- \'Ecole Nationale Sup\'erieure de l'\'Electronique et de ses Applications} \hfill \textit{2025 -- 2027} \\
    \small{Double dipl\^ome d'Ing\'enieur -- Cybersecurity, Data Science \& Intelligence Artificielle.}
    \item \textbf{\'Ecole Polytechnique de Yaound\'e (1\textsuperscript{\`ere} \'ecole d'ing\'enieur CEMAC)} \hfill \textit{2021 -- 2025} \\
    \small{Dipl\^ome d'Ing\'enieur de Conception (Master 2) : \textbf{Math\'ematiques Appliqu\'ees}, Algorithmique, G\'enie Logiciel.}
\end{itemize}

\section{Comp\'etences Techniques}
\begin{itemize}[leftmargin=*, nosep]
    \item \small{\textbf{IA Generative \& LLM :} LangChain, Gemini, GPT, RAG avanc\'ee, prompt engineering, agents multi-etapes, fine-tuning de LLMs, evaluation et evals automatis\'ees de quality.}
    \item \small{\textbf{NLP \& Transformers :} BERT, GPT-2, fine-tuning, classification de textes, extraction d'entites, analyse de sentiment, embeddings contextuels, Hugging Face ecosysteme.}
    \item \small{\textbf{Data Science \& ML :} XGBoost, LightGBM, ensemble methods, hyperparameter tuning, feature engineering avancee, cross-validation, A/B testing, scoring et probabilites.}
    \item \small{\textbf{Data Engineering \& Python :} ETL pipelines, pandas, PySpark, SQL (PostgreSQL), APIs REST, Flask, Docker, versioning (Git), monitoring et logging.}
    \item \small{\textbf{Langues :} Anglais courant (TOEFL approx. 100), Francais natif, Camerounais (Yemba).}
\end{itemize}

\section{Exp\'eriences \& Projets}
\begin{itemize}[leftmargin=*, label={-}]

\item \textbf{Stage INRIA Saclay -- \'Equipe TRIBE (Confiance \& Transparence)} \hfill \textit{Avril-Ao\^ut 2026} \\
\textit{Palaiseau (IP Paris)} \hfill \textit{Python, BERT, Transformers, NLP, scikit-learn, pandas}
\vspace{-2pt}
    \begin{itemize}[leftmargin=0.4cm, nosep, label=\tiny$\bullet$]
        \item \small{Pipeline NLP complet pour analyse donnees financieres/mobilite : fine-tuning BERT sur corpus de politiques confidentialite, tokenization specialisee, classification multi-labels pour risques RGPD.}
        \item \small{Score de confiance interpreatable (0-100) agrageant risques de fuite donnees, non-conformites et qualite de transparence. Backend Python optimise pour latence production < 100ms par requete.}
        \item \small{Analyse statique et NLP combinaison pour identification SDKs/APIs posant risques. Pipeline complet extraction features, scoring probabiliste, integrabilite dans systemes decision bancaires.}
        \item \small{Collaboration LVMT/Ecole Ponts, programme PEPR MOBIDEC (financement CNRS/ANR). Redaction scientifique anglais, communication resultats aupres équipe de recherche.}
    \end{itemize}
\vspace{3pt}

\item \textbf{Assistant Recherche MILA (Montreal Institute Learning Algorithms)} \hfill \textit{Juin-Sept 2025} \\
\textit{Remote -- Phenomene Grokking et Generalisation} \hfill \textit{PyTorch, TensorFlow, research methods, mathematics}
\vspace{-2pt}
    \begin{itemize}[leftmargin=0.4cm, nosep, label=\tiny$\bullet$]
        \item \small{Etude phénomène grokking (generalisation tardive extreme) : reproduction experiences SOTA sur MLPs, implementation complète taches algebriques (modular arithmetic, matrix inversion).}
        \item \small{Tests statistiques rigoureux, validation croisee, analyse convergence dynamiques apprentissage. Hypotheses mathematiques sur mecanismes formation features et double descent.}
        \item \small{Visualisations manifold learning (t-SNE, PCA), interpretabilite representations internes. Resultats reproductibles sur GitHub avec documentation scientifique detaillee.}
    \end{itemize}
\vspace{3pt}

\item \textbf{Systeme Multi-Agents Agricole avec Google ADK} \hfill \textit{2024-2025} \\
\textit{RAG, Orchestration Agents, APIs Externes} \hfill \textit{Google ADK, LangChain, Gemini, RAG, FastAPI, Docker}
\vspace{-2pt}
    \begin{itemize}[leftmargin=0.4cm, nosep, label=\tiny$\bullet$]
        \item \small{4 agents IA collaboratifs (Planning, Meteo, Marche, Recommandations) orches-tres par Google ADK avec LLM Gemini base, resolution conflicts multi-agents, planning hierarchique.}
        \item \small{RAG enrichissement contextuel : integration APIs meteo (OpenWeather), donnees marches agricoles, recommandations cultures multi-criteres (rendement, rentabilite, climat). Vector stores pour optimisation retrieval.}
        \item \small{Architecture scalable FastAPI + Docker, logs structurees pour traceabilite decisions agents. Evaluation agents sur benchmarks agronomiques, rapport impact utilisateur final.}
    \end{itemize}
\vspace{3pt}

\item \textbf{UROP 2024 -- Classification ECG depuis Photographies} \hfill \textit{2024} \\
\textit{Computer Vision Medicale} \hfill \textit{TensorFlow, Keras, Computer Vision, signal processing}
\vspace{-2pt}
    \begin{itemize}[leftmargin=0.4cm, nosep, label=\tiny$\bullet$]
        \item \small{Modele CNN detection/classification anomalies cardiaques from ECG photographe (approche end-to-end). Feature extraction signaux ECG (onde P, QRS, T), filtrage adaptatif bruit.}
        \item \small{Metriques evaluation cliniques : sensibilite, specificite, AUC-ROC. Validation sur dataset cardiovasculaire, collaboration mentors Google DeepMind pour approches SOTA.}
    \end{itemize}
\vspace{3pt}

\item \textbf{Urban Waste Management -- Optimisation Collecte Dechets} \hfill \textit{2024} \\
\textit{1\textsuperscript{er} Prix CITS National contre 75 equipes} \hfill \textit{ML, XGBoost, algorithmes optimisation, IoT}
\vspace{-2pt}
    \begin{itemize}[leftmargin=0.4cm, nosep, label=\tiny$\bullet$]
        \item \small{Prediction besoins collecte via XGBoost sur IoT sensor data (temperature, humidite, jour). TSP solver avec 2-opt, 3-opt pour optimisation routes : reduction 20\% couts operationnels.}
        \item \small{Proposition paper de recherche, publication en examination dans conference Smart Cities africaines. Transfert technologie vers gouvernements locaux.}
    \end{itemize}
\vspace{3pt}

\item \textbf{IndabaX Africa 2025 -- 1\textsuperscript{er} Prix (IA Sante)} \hfill \textit{2025} \\
\textit{Data Cleaning \& Streamlit Dashboard} \hfill \textit{Python, Pandas, Flask, Streamlit, SQL}
\vspace{-2pt}
    \begin{itemize}[leftmargin=0.4cm, nosep, label=\tiny$\bullet$]
        \item \small{Pipeline nettoyage complet donnees medicales (outlier detection, imputation missing values, normalization robuste sur dataset multimodal). API Flask avec endpoints RESTful authentifies.}
        \item \small{Dashboard Streamlit interactif : analyse exploratoire, visualisation balances, export rapports. Interface accessible, deploiement cloud-ready avec monitoring.}
    \end{itemize}
\vspace{3pt}

\end{itemize}

\section{Distinctions \& Leadership}
\begin{itemize}[leftmargin=*, label=\tiny$\bullet$, nosep]
    \item \small{\textbf{1\textsuperscript{er} Prix IndabaX Africa 2025}} (IA Sante) -- meilleur projet data science et ML santé sur continent africain, reconnaissance excellence technique.}
    \item \small{\textbf{1\textsuperscript{er} Prix National CITS 2024}} (Urban Waste Management) -- competences optimisation algorithmique et AI for Good.}
    \item \small{\textbf{Stage INRIA Saclay}} (Avril-Aout 2026) -- selection competitive, equipe TRIBE, collaboration CNRS/ANR programme PEPR MOBIDEC.}
    \item \small{\textbf{Certifications en Cours}} : Google Cloud Associate Data Engineer (juin 2026), AWS ML Specialty (septembre 2026).}
\end{itemize}

\end{document}
